\section{技术可行性论证}\label{sec:feasibility}

\subsection{神经算子理论保证}

E3A框架的核心算法性能预测模块基于DeepONet，其数学基础为Chen \& Chen（1995）关于神经算子的万能逼近定理\cite{chen1995}：

\begin{tcolorbox}[hypothesisstyle,title={定理：神经算子万能逼近性（Chen \& Chen, 1995）}]
设$\calF: C(K_1;\R^d) \to C(K_2;\R)$为定义在紧集上的连续非线性算子，$\sigma$为非多项式连续激活函数。则对任意$\varepsilon > 0$，存在正整数$p$与神经网络参数，使得
$$
\sup_{u \in C(K_1;\R^d)} \left|\calF(u)(y) - \sum_{i=1}^{p} b_i(u) \cdot t_i(y)\right| < \varepsilon, \quad \forall y \in K_2.
$$
\end{tcolorbox}

在ECC性能预测中，$u$对应信道三元组$\calT$（连续函数），$y$对应SNR评估点，$\calF$对应BER/FER性能算子。定理保证框架在理论上具备无限逼近精度，与码型和信道类型无关。

\subsection{DeepONet在通信系统中的先验验证}

Lu等（2021）\cite{lu2021}在\textit{Nature Machine Intelligence}上首次报告DeepONet在物理场算子学习中的实验验证，证明神经算子在复杂非线性算子逼近中兼具高精度（误差$<1\%$）与高效推理（加速比$>10^4\times$）。本框架将此范式首次系统化迁移至ECC性能预测，理论基础充分，先前方法验证充分。

\subsection{已有相关工作支撑}

\begin{table}[htbp]
\centering
\caption{E3A各技术模块的先前工作支撑}
\label{tab:prior_work}
\setlength{\tabcolsep}{4pt}
\rowcolors{2}{lightblue!15}{white}
\begin{tabular}{@{}L{2.5cm}L{4.5cm}L{4.5cm}L{2.0cm}@{}}
\toprule
\textbf{技术模块} & \textbf{代表性先前工作} & \textbf{关键进展} & \textbf{本工作突破} \\
\midrule
因子图GNN & Nachmani等（2018）\cite{nachmani2018} & GNN显著优于BP & 推理不译码 \\
图变换器 & Dwivedi \& Bresson（2021）\cite{dwivedi2021} & 图注意力+位置编码 & 迁移至ECC图 \\
神经算子 & Lu等（2021）\cite{lu2021} & DeepONet理论与实验 & 首次用于ECC \\
硬件代理模型 & Ustun等（2020）\cite{ustun2020} & HLS性能机器学习预测 & 扩展至ECC \\
LLM for EDA & Chang等（2024）\cite{chang2024} & LLM辅助RTL生成 & 形式约束解析 \\
\bottomrule
\end{tabular}
\end{table}

\subsection{计算量估算}

\textbf{训练阶段}：性能预测代理模型单次训练约需$10^5$个（码参数，信道参数，性能标签）样本，可在$4\times$A100 GPU上于24小时内完成；硬件代理模型训练样本通过Design Compiler批量综合产生，单批次（100配置）约需8小时，并行化后可于1周内建库。\textbf{推理阶段}：单次神经算子推理耗时约$0.5$~ms（A100）或约$5$~ms（CPU），满足交互式设计需求。\textbf{与蒙特卡洛基准对比}：蒙特卡洛仿真达到BER$=10^{-9}$置信度需$\sim10^{11}$次仿真，在现代服务器上约需3--10小时；代理模型将其压缩至毫秒级，\textbf{加速比约$10^6\sim10^7\times$}。